% META: {"id": "TNSI-ARCH-EVAL-B-corrige", "chapitre": "TNSI-ARCHITECTURES-MATERIELLES-SY", "type_objet": "correction", "evaluation_ref": "TNSI-ARCH-EVAL-B", "capacites": ["T-ARCH-03", "T-ARCH-04A", "T-ARCH-04B"], "mode_creation": "ex_nihilo", "sources_inspiration": [], "status": "approved"}
% BEGIN-VERIFY
% def bfs_sauts(graphe_couts, depart, arrivee):
%     visites = {depart: 0}
%     file = [depart]
%     while file:
%         courant = file.pop(0)
%         if courant == arrivee:
%             return visites[courant]
%         for voisin in graphe_couts[courant]:
%             if voisin not in visites:
%                 visites[voisin] = visites[courant] + 1
%                 file.append(voisin)
%     return None
% def plus_court_cout(graphe_couts, depart, arrivee):
%     couts = {depart: 0}
%     a_traiter = [depart]
%     while a_traiter:
%         courant = min(a_traiter, key=lambda s: couts[s])
%         a_traiter.remove(courant)
%         if courant == arrivee:
%             break
%         for voisin, poids in graphe_couts[courant].items():
%             nouveau_cout = couts[courant] + poids
%             if voisin not in couts or nouveau_cout < couts[voisin]:
%                 couts[voisin] = nouveau_cout
%                 if voisin not in a_traiter:
%                     a_traiter.append(voisin)
%     return couts[arrivee]
% reseau = {
%     "A": {"B": 1, "D": 8},
%     "B": {"A": 1, "C": 1},
%     "C": {"B": 1, "D": 1},
%     "D": {"A": 8, "C": 1},
% }
% assert bfs_sauts(reseau, "A", "D") == 1
% assert plus_court_cout(reseau, "A", "D") == 3
% END-VERIFY
\begin{corrige}{TNSI-ARCH-EVAL-B-EX1}
\textbf{Q1.} RIP choisirait la route directe \texttt{A-D} (un seul saut), quel que soit son coût élevé. \lstinline{bfs_sauts(reseau, "A", "D")} renvoie $1$.

\textbf{Q2.} \lstinline{plus_court_cout(reseau, "A", "D")} renvoie $3$ : la route la moins coûteuse est \texttt{A-B-C-D} (coût $1+1+1=3$), bien moins chère que la route directe (coût $8$). Elle \textbf{ne coïncide pas} avec la route RIP : RIP privilégie le nombre de sauts, OSPF le coût cumulé.
\end{corrige}

\begin{corrige}{TNSI-ARCH-EVAL-B-EX2}
\textbf{Q1.} Le chiffrement symétrique utilise une \textbf{unique clé secrète} partagée pour chiffrer et déchiffrer ; le chiffrement asymétrique utilise une \textbf{paire de clés} (publique/privée) : ce qui est chiffré avec la clé publique ne se déchiffre qu'avec la clé privée correspondante.

\textbf{Q2.} Le chiffrement asymétrique est beaucoup plus lent à calculer que le chiffrement symétrique. L'utiliser pour \textbf{tout} l'échange (pages web, formulaires, fichiers) ralentirait considérablement la connexion. HTTPS l'utilise donc seulement pour échanger, en toute sécurité, une clé symétrique de session au début de la connexion ; le reste des données est ensuite chiffré avec cette clé symétrique, bien plus rapide.
\end{corrige}
